
\section{Related Work}
\label{sec:related-work}

We connect three strands of research: the reception of minority judgments,
the moral reasons exchanged in online discussions, and interventions that
shape readers' exposure to different views. Together, they motivate
examining dissent through community responses, the content of replies,
and readers' consideration of alternative reasons.

\subsubsection*{Minority Reception}

Research on minority influence distinguishes public acceptance from
effects on private judgments \cite{wood1994minority}. Dissenters can
attract communication while also facing rejection
\cite{schachter1951deviation}. These findings suggest that attention
to a minority position need not indicate approval of it.

Online feedback similarly separates evaluation from engagement.
Prior ratings can influence subsequent evaluations
\cite{muchnik2013social}, and Reddit voting reflects local norms
and practices \cite{graham2021sociomateriality}. Negatively rated
comments can nevertheless attract more replies and larger reply
trees \cite{davis2021emotional}. In AITA, the relative valuation
of agreement and dissent also varies with the strength of the
developing consensus \cite{mellody2025consensus}. Together, these
studies motivate examining minority reception through both scores
and conversational engagement, while accounting for the discussion's
level of agreement.

Reception may also differ with the moral direction of dissent.
Public blame can itself invite scrutiny and demands for justification
\cite{malle2022socially}. Blaming dissent calls for condemnation
against a defending majority, whereas defending dissent resists
the majority's blame. Both oppose the prevailing judgment, but make
different moral claims. In \textbf{RQ1}, we examine minority reception
across high- and low-agreement discussions and compare blaming and
defending dissent within high-agreement discussions, using majority
comments as a reference for scores, reply counts, and reply depth.


\subsubsection*{Minority Discourse and Moral Reasoning}

Research on AITA characterizes everyday moral dilemmas and the
situational and relational contexts of disagreement
\cite{nguyen2022mapping,shi2025moral}. Within these discussions,
judgments are both expressions of a moral position and contributions
to exchanges of reasons. Understanding minority discourse therefore
requires examining how firmly judgments are expressed, how others
respond, and what moral reasons emerge.

%Anticipated evaluation can shape expression. Accountability may
%encourage accommodation to an audience or bolstering of prior
%commitments \cite{tetlock1989strategies}. 
When people expect others to judge their views, they may try to
please their audience or defend what they already believe
\cite{tetlock1989strategies}.
Research on the spiral
of silence links willingness to express minority views to opinion
climate and attitude certainty \cite{matthes2010spiral}, while
anticipated personal attacks are also associated with willingness
to speak \cite{neubaum2018sanctions}. Expressions of doubt receive
fewer recommendations in online news discussions
\cite{evans2023doubt}. These findings motivate examining expressed
certainty, while distinguishing linguistic confidence from private
belief and recognizing that observed patterns may reflect selective
participation.

The reasons supporting judgments provide a complementary dimension
of discourse. Computational research represents situation-specific
moral norms \cite{forbes2020social} and potentially conflicting
values, rights, and duties \cite{sorensen2024value}.
\citet{xi2024morality} extract and categorize rationale elements
in AITA comments. These approaches characterize moral content
within contributions; examining replies additionally reveals
whether exchanges introduce reasons absent from the parent comment.

Dissent can encourage divergent thinking \cite{nemeth1986differential},
and deliberative research distinguishes agreement from the ability
to articulate reasons for different positions.
\emph{Argument repertoire} encompasses reasons supporting one's own
and opposing views \cite{cappella2002argument}, with exposure to
disagreement associated with a greater ability to offer opposing
reasons \cite{price2002disagreement}. Building on these distinctions,
\textbf{RQ2} examines expressed certainty, reply orientation, and
moral reasons in comments and replies. Measuring support and
challenge separately from additional reasons allows both kinds
of replies to extend a parent's reasoning.

\subsubsection*{Information Exposure and Reader Responses}

Research on information exposure distinguishes available content,
algorithmic selection, and users' consumption choices
\cite{bakshy2015exposure}. Much intervention research concerns
political news and partisan attitudes, with varying effects of
exposure to opposing or like-minded sources
\cite{bail2018exposure,levy2021social,nyhan2023likeminded}.
These findings motivate specifying both what an intervention
surfaces and which reader outcomes it addresses.

Interventions differ in how they shape exposure. Social Mirror
supports reflection through network visualization and account
recommendations \cite{gillani2018mirror}, while Opinion Space
and ConsiderIt introduce dedicated interfaces for navigating
perspectives or reasons
\cite{faridani2010opinion,kriplean2012considerit}. Other approaches
change which existing contributions appear prominently, through
comment diversification \cite{northcutt2017diversification} or
feed reranking \cite{piccardi2025reranking}. Thus, changing
selection within familiar displays offers an alternative to
introducing additional representations.

Closer experimental precedents compare agreeing and disagreeing
comment pairs in Facebook-style displays
\cite{dutceac2022cueing}. Moral-judgment studies also examine
exposure to peer evaluations and arguments
\cite{kelly2017moral}, or aggregate controversy and group-certainty
signals derived from AITA discussions
\cite{chen2026controversy}. These studies establish connections
between social information and reader responses, while leaving
room to examine how selecting existing minority threads affects
consideration of moral reasons.

In \textbf{RQ3}, we test a selection intervention within a
Reddit-style comment display: the highest-scoring minority
thread replaces the second majority thread, retaining original
wording, scores, and accompanying replies. Minority status is
defined relative to the discussion, rather than the reader's
prior position. The design surfaces competing judgments about
the same interpersonal dilemma without adding visualizations
or aggregate judgment summaries. We examine readers' judgment
certainty and reported consideration of alternative moral
reasons, alongside agreement and perceived comment influence,
to assess responses beyond adoption of the minority verdict.